\section*{Abstract}
\phantomsection
\addcontentsline{toc}{section}{Abstract}

\textbf{Background.}
Pancreatic neuroendocrine tumours are rare, and liver metastasis is associated
with more complex diagnosis-time risk assessment and prognosis. This study
examined whether published diagnostic and prognostic workflows could be
reconstructed using an available Surveillance, Epidemiology, and End Results
(SEER) cohort. It also evaluated whether tabular foundation models,
Tabular Prior-Data Fitted Network (TabPFN) representations, multilayer
perceptrons (MLPs) or synthetic-data augmentation improved prediction.

\textbf{Methods.}
The November 2023 SEER release yielded 2,672 eligible diagnostic records,
including 462 patients with liver metastasis at diagnosis; those 462 patients
formed the prognostic cohort. The reference workflows were reconstructed using
the available variables, and the resulting models were applied without
modification to an independent University Hospital Erlangen cohort of 30
patients, including 11 with liver metastasis. Supervised TabPFN embeddings and
complete synthetic cohorts containing predictors and outcomes were evaluated
for predictive information, structural similarity, category support and
preservation of clinical relationships. Conventional models, tabular
foundation models and prespecified MLPs were compared using the same repeated
SEER training and test partitions. Synthetic augmentation added 1,000, 5,000,
10,000 or 15,000 generated records to each real training partition, while
performance was measured only in held-out real patients.

\textbf{Results.}
The reference workflows recovered strong diagnostic discrimination and
meaningful prognostic risk ordering, although the unavailable original SEER
release prevented exact replication. In the independent Erlangen cohort, the
primary diagnostic nomogram achieved an area under the receiver operating
characteristic curve (ROC AUC) of 0.934 (95\% confidence interval
0.809--1.000; $n=29$), and the primary prognostic nomogram achieved a Harrell
concordance index of 0.833 (0.625--1.000; $n=11$). Both estimates were
imprecise because of the small sample. In the repeated SEER evaluation,
TabPFN-3, TabICLv2 and TabDPT achieved diagnostic ROC AUCs of 0.902--0.903,
with improvements of 0.0074--0.0088 over logistic regression on the same
patients. The clinical MLP showed no general advantage. No foundation model or
MLP showed a reliable prognostic concordance gain. Random survival forest had
the lowest integrated Brier score (0.153), improving on Cox elastic net by
0.0071. At the largest augmentation size, diagnostic XGBoost improved ROC AUC
by 0.0126 and reduced Brier score by 0.0060; other diagnostic models did not
improve consistently, and prognostic augmentation provided no benefit.
Embeddings retained predictive information and local-neighbourhood structure,
whereas the synthetic cohorts lost grade-4 support and shifted survival and
several relationships between variables.

\textbf{Conclusions.}
The benefits of foundation models and synthetic data were specific to the
task, model and evaluation design. The strongest evidence supports small
internal diagnostic gains and a selective augmentation benefit for XGBoost,
not a general prognostic or clinical advantage. The external results are
encouraging for the reproduced reference models but remain too small to
establish subgroup safety or clinical usefulness. Larger independent
validation is required before any model can support clinical decisions.

\vfill
\noindent\textbf{Keywords:} pancreatic neuroendocrine tumour; SEER; tabular
foundation models; TabPFN; synthetic data; survival analysis; clinical prediction
